% Part: sets-functions-relations
% Chapter: infinite
% Section: dedekind
%
\documentclass[../../../include/open-logic-section]{subfiles}

\begin{document}

\olfileid{sfr}{infinite}{dedekind}
\olsection{Dedekindalgebra's}

De natuurlijke getallen moeten niet alleen oneindig zijn. Ze moeten ook
bepaalde (algebraïsche) eigenschappen hebben, zodat optellen,
vermenigvuldigen enzovoort goed werken.

Dedekind nam de \emph{opvolgerfunctie} als uitgangspunt. Die functie
stuurt elk getal naar zijn opvolger. Daarna beschreef hij de getallen
met deze drie eigenschappen:
\begin{enumerate}
	\item Er is een getal, $0$, dat niet de opvolger van een getal is
	\\anders gezegd: $0 \notin \ran{s}$
	\\in symbolen: $\forall x\ s(x) \neq 0$
	\item Verschillende getallen hebben verschillende opvolgers
	\\dus is $s$ injectief: verschillende invoeren krijgen verschillende uitkomsten
	\\in symbolen: $\forall x \forall y (s(x) = s(y) \lif x = y)$
	\item\ollabel{repeatedapplication} Elk getal krijg je door bij $0$ te
	beginnen en steeds de opvolgerfunctie toe te passen.
\end{enumerate}
De eerste twee eigenschappen kunnen we makkelijk in eerste-ordelogica
opschrijven (zie hierboven). Maar voor \olref{repeatedapplication}
hebben we meer nodig dan alleen eerste-ordelogica. Dedekind wist deze
eigenschap, \olref{repeatedapplication}, als volgt met verzamelingen te
omschrijven:
\begin{enumerate}
	\item[3$'$.] De natuurlijke getallen vormen de kleinste verzameling
	die \emph{gesloten is onder de opvolgerfunctie}. Dat betekent: pas je
	$s$ toe op een willekeurig element van de verzameling, dan krijg je
	opnieuw een element van die verzameling.
\end{enumerate}
We moeten dit nu precies uitschrijven.

\begin{defn}\ollabel{Closure}
	Neem een functie $f$. We noemen een verzameling $X$ $f$-\emph{gesloten}
	{precies als} $(\forall x \in X)f(x) \in X$. De functie brengt
	elementen van die verzameling dan niet buiten de verzameling. Voor
	elke $o$ leggen we nu vast:
	$$\closureofunder{f}{o} = \bigcap\Setabs{X}{o \in X\text{ en }X
\text{ $f$-gesloten is}}$$
\end{defn}

Dus $\closureofunder{f}{o}$ is de doorsnede van alle $f$-gesloten
verzamelingen waar $o$ een element van is: het deel dat al die
verzamelingen gemeen hebben. Daarom is $\closureofunder{f}{o}$
intuïtief de \emph{kleinste} $f$-gesloten verzameling waar $o$ een
element van is. Het volgende resultaat maakt dat precies:
\begin{lem}\ollabel{closureproperties}
	Voor elke verzameling $A$, elke functie $f \colon A \to A$ en elke
	$o \in A$:
	\begin{enumerate}
		\item\ollabel{closurehaselem} $o \in \closureofunder{f}{o}$; en
		\item\ollabel{closureclosed} $\closureofunder{f}{o}$ is $f$-gesloten; en
		\item\ollabel{closuresmallest} als $X$ $f$-gesloten is en $o \in
		X$, dan $\closureofunder{f}{o} \subseteq X$
	\end{enumerate}
\end{lem}

\begin{proof}
Er is minstens één $f$-gesloten verzameling waar $o$ een element van
is, namelijk $\ran{f}\cup \{o\}$. Daarom bestaat
$\closureofunder{f}{o}$, de doorsnede van \emph{alle} zulke
verzamelingen. We moeten nu
\olref{closurehaselem}--\olref{closuresmallest} nalopen.

Voor \olref{closurehaselem}: $o \in \closureofunder{f}{o}$. Alle
verzamelingen in de doorsnede hebben $o$ als element, dus hun doorsnede
ook.

Voor \olref{closureclosed}: neem $x \in \closureofunder{f}{o}$. Als
$o \in X$ en $X$ $f$-gesloten is, dan zit $x \in X$ en zit ook
$f(x) \in X$, want $X$ is $f$-gesloten. Dit geldt voor elke verzameling
in de doorsnede, dus $f(x) \in \closureofunder{f}{o}$.

Voor \olref{closuresmallest} gebruiken we een algemene regel: als
$X \in C$, dan $\bigcap C \subseteq X$.
\end{proof}

Nu kunnen we de volgende definitie geven:

\begin{defn}
Een \emph{Dedekindalgebra} bestaat uit drie dingen: een verzameling
$A$, een functie $f \colon A \to A$ en een gekozen element $o \in A$.
Ze moeten aan deze voorwaarden voldoen:
	\begin{enumerate}
		\item \ollabel{ded:proper} $o \notin \ran{f}$
		\item \ollabel{ded:injection} $f$ is injectief
		\item \ollabel{ded:closure} $A = \closureofunder{f}{o}$
	\end{enumerate}
\end{defn}

Omdat $A = \closureofunder{f}{o}$, weten we uit het vorige resultaat
dat $A$ de kleinste $f$-gesloten verzameling is waar $o$ een element
van is. Elke Dedekindalgebra is Dedekind-oneindig. Dat zie je direct aan
voorwaarden \olref{ded:proper} en \olref{ded:injection} van de
definitie. Maar er geldt ook iets in de omgekeerde richting: van elke
Dedekind-oneindige verzameling kunnen we een Dedekindalgebra maken.

\begin{thm}\ollabel{thm:DedekindInfiniteAlgebra}
Als er een Dedekind-oneindige verzameling bestaat, dan bestaat er een
Dedekindalgebra.
\end{thm}

\begin{proof}
Neem een Dedekind-oneindige verzameling $D$. Er is dan een injectie $g
\colon D \to D$ en een element $o \in D \setminus \ran{g}$. Neem nu
$A = \closureofunder{g}{o}$. Volgens \olref{closureproperties} bestaat
$A$ en geldt $o \in A$. Neem $f = \funrestrictionto{g}{A}$. We laten
zien dat $A, f, o$ samen een Dedekindalgebra vormen.

Voor \olref{ded:proper}: $o \notin \ran{g}$ en $\ran{f} \subseteq
\ran{g}$, dus $o\notin \ran{f}$.

Voor \olref{ded:injection}: $g$ is een injectie op $D$ en $f
\subseteq g$, dus is de beperking ook injectief.

Voor \olref{ded:closure}: volgens \olref{closureproperties} is $A$
$g$-gesloten; dus is $A$ ook $f$-gesloten. Daarom is
$\closureofunder{f}{o} \subseteq A$ volgens
\olref{closureproperties}. Verder is $\closureofunder{f}{o}$ zelf
$f$-gesloten en geldt $f = \funrestrictionto{g}{A}$. Daaruit volgt dat
$\closureofunder{f}{o}$ ook $g$-gesloten is. Dus
$A \subseteq \closureofunder{f}{o}$ volgens
\olref{closureproperties}.
\end{proof}

\end{document}